%! AP Statistics — Unit 9, Lesson 9.2 — Exit Ticket (Answer Key)
\documentclass[10pt]{article}
\usepackage{apstats-article}
\usepackage{apstats-key}

\begin{document}

\pageheader{Unit 9, Lesson 9.2}{Exit Ticket --- Answer Key}
\namedateperiod

\vspace{0.1in}
A nutritionist randomly selects 20 adults and records their daily calorie intake ($x$,
in hundreds of calories) and resting metabolic rate ($y$, calories/day).
Computer output gives $b = 18.4$, $SE_b = 4.2$, and $n = 20$.
Assume all LINER conditions are met. Use $t^* = 2.101$ for a 95\% CI ($df = 18$).

\vspace{0.15in}
\begin{enumerate}

    \item Construct a 95\% confidence interval for the true slope $\beta$. Show your work.

          \vspace{0.05in}
          $b \pm t^* \cdot SE_b
          = $ \ans{$18.4$} $\pm$ \ans{$2.101$} $\times$ \ans{$4.2$}
          $= $ \ans{$18.4 \pm 8.824$}
          $= ($ \ans{$9.576$},\; \ans{$27.224$}$)$

          \vspace{0.1in}

    \item Interpret this confidence interval in context.

          \vspace{0.05in}
          \ans{We are 95\% confident that the true slope of the population regression
          line relating resting metabolic rate to daily calorie intake for adults is
          between approximately 9.6 and 27.2 calories per day per hundred calories consumed.}

          \vspace{0.1in}

    \item Does this CI provide evidence that $\beta \ne 0$? Justify your answer.

          \vspace{0.05in}
          \ans{Yes. Since 0 is not contained in the interval $(9.6,\ 27.2)$, this provides
          evidence that $\beta \ne 0$; there is a positive linear relationship between
          daily calorie intake and resting metabolic rate.}

\end{enumerate}

\end{document}
